% Options for packages loaded elsewhere
\PassOptionsToPackage{unicode}{hyperref}
\PassOptionsToPackage{hyphens}{url}
%
\documentclass[
  journal,
  onecolumn]{IEEEtran}
\usepackage{amsmath,amssymb}
\usepackage{lmodern}
\usepackage{iftex}
\ifPDFTeX
  \usepackage[T1]{fontenc}
  \usepackage[utf8]{inputenc}
  \usepackage{textcomp} % provide euro and other symbols
\else % if luatex or xetex
  \usepackage{unicode-math}
  \defaultfontfeatures{Scale=MatchLowercase}
  \defaultfontfeatures[\rmfamily]{Ligatures=TeX,Scale=1}
\fi
% Use upquote if available, for straight quotes in verbatim environments
\IfFileExists{upquote.sty}{\usepackage{upquote}}{}
\IfFileExists{microtype.sty}{% use microtype if available
  \usepackage[]{microtype}
  \UseMicrotypeSet[protrusion]{basicmath} % disable protrusion for tt fonts
}{}
\makeatletter
\@ifundefined{KOMAClassName}{% if non-KOMA class
  \IfFileExists{parskip.sty}{%
    \usepackage{parskip}
  }{% else
    \setlength{\parindent}{0pt}
    \setlength{\parskip}{6pt plus 2pt minus 1pt}}
}{% if KOMA class
  \KOMAoptions{parskip=half}}
\makeatother
\usepackage{xcolor}
\setlength{\emergencystretch}{3em} % prevent overfull lines
\providecommand{\tightlist}{%
  \setlength{\itemsep}{0pt}\setlength{\parskip}{0pt}}
\setcounter{secnumdepth}{-\maxdimen} % remove section numbering
\newlength{\cslhangindent}
\setlength{\cslhangindent}{1.5em}
\newlength{\csllabelwidth}
\setlength{\csllabelwidth}{3em}
\newlength{\cslentryspacingunit} % times entry-spacing
\setlength{\cslentryspacingunit}{\parskip}
\newenvironment{CSLReferences}[2] % #1 hanging-ident, #2 entry spacing
 {% don't indent paragraphs
  \setlength{\parindent}{0pt}
  % turn on hanging indent if param 1 is 1
  \ifodd #1
  \let\oldpar\par
  \def\par{\hangindent=\cslhangindent\oldpar}
  \fi
  % set entry spacing
  \setlength{\parskip}{#2\cslentryspacingunit}
 }%
 {}
% Add lowercase alias for newer pandoc versions (no parameters)
\newenvironment{cslreferences}
 {% don't indent paragraphs
  \setlength{\parindent}{0pt}
  % turn on hanging indent
  \let\oldpar\par
  \def\par{\hangindent=\cslhangindent\oldpar}
  % set entry spacing
  \setlength{\parskip}{\cslentryspacingunit}
 }%
 {}
\usepackage{calc}
\newcommand{\CSLBlock}[1]{#1\hfill\break}
\newcommand{\CSLLeftMargin}[1]{\parbox[t]{\csllabelwidth}{#1}}
\newcommand{\CSLRightInline}[1]{\parbox[t]{\linewidth - \csllabelwidth}{#1}\break}
\newcommand{\CSLIndent}[1]{\hspace{\cslhangindent}#1}
\usepackage[margin=1in]{geometry}
\usepackage{longtable}
\usepackage{amsmath}
\usepackage{amsfonts}
\usepackage{array}
\usepackage{booktabs}
\usepackage{caption}
\renewcommand{\thetable}{S\arabic{table}}
\renewcommand{\thefigure}{S\arabic{figure}}
\renewcommand{\thesection}{S\arabic{section}}
\ifLuaTeX
  \usepackage{selnolig}  % disable illegal ligatures
\fi
\IfFileExists{bookmark.sty}{\usepackage{bookmark}}{\usepackage{hyperref}}
\IfFileExists{xurl.sty}{\usepackage{xurl}}{} % add URL line breaks if available
\urlstyle{same} % disable monospaced font for URLs
\hypersetup{
  pdftitle={Supplementary Material: Lite Learning: Efficient Crop Classification in Tanzania Using Feature Extraction with Machine Learning \& Crowd Sourcing},
  hidelinks,
  pdfcreator={LaTeX via pandoc}}

\title{Supplementary Material: Lite Learning: Efficient Crop
Classification in Tanzania Using Feature Extraction with Machine
Learning \& Crowd Sourcing}
\author{
Michael L. Mann\thanks{The George Washington University, Washington DC
20052, Corresponding author. Email mmann1123@gmail.com} \and 
Lisa Colson\thanks{USDA Foreign Agricultural Service, Washington DC
20250} \and 
Rory Nealon\thanks{USAID GeoCenter, Washington DC 20523} \and 
Ryan Engstrom\thanks{The George Washington University, Washington DC
20052} \and 
Stellamaris Nakacwa\thanks{YouthMappers, Texas Tech University, Lubbock
TX 79409}}
\date{}

\begin{document}
\maketitle

\hypertarget{s-i.-time-series-feature-definitions}{%
\section{S-I. Time Series Feature
Definitions}\label{s-i.-time-series-feature-definitions}}

This document is the online supplement to the main article \emph{``Lite
Learning: Efficient Crop Classification in Tanzania Using Feature
Extraction with Machine Learning \& Crowd Sourcing.''} It provides the
full list of time-series features referenced in the main text (Methods,
``Time Series Features'').

The features were extracted on a per-pixel basis from the 2023
Sentinel-2 monthly composites (bands B2, B6, B11, B12, the Enhanced
Vegetation Index (EVI), and hue) using the \texttt{xr\_fresh} toolkit
(Mann, Michael L. 2024). They summarize the temporal dynamics of crop
growth and development across the growing season, spanning basic
statistical descriptors, changes over time, and distribution-based
metrics. Table S1 lists each statistic, a short description, and its
defining equation. Throughout, \(x_i\) denotes the value of the band or
index at time step \(i\) in a pixel's monthly time series of length
\(n\), \(\bar{x}\) its mean, \(\sigma\) its standard deviation,
\(\mu_4\) its fourth central moment, and \(l\) an autocorrelation lag.

\renewcommand{\arraystretch}{1.5}
\begin{longtable}{|p{4cm}|p{5cm}|p{6cm}|}
\caption{Time-series features extracted with \texttt{xr\_fresh}, with descriptions and defining equations.}\label{tab:ts_features}\\
\hline
\textbf{Statistic} & \textbf{Description} & \textbf{Equation} \\
\hline
\endhead
Absolute energy &  sum over the squared values & $E = \sum_{i=1}^n x_i^2$ \\
Absolute Sum of Changes  & sum over the absolute value of consecutive changes in the series  & $ \sum_{i=1}^{n-1} \mid x_{i+1}- x_i \mid $ \\
Autocorrelation (1, 2 \& 3 month lag) & Correlation between the time series and its lagged values & $\frac{1}{(n-l)\sigma^{2}} \sum_{t=1}^{n-l}(x_{t}-\bar{x})(x_{t+l}-\bar{x})$\\
Count Above Mean & Number of values above the mean & $N_{\text{above}} = \sum_{i=1}^n (x_i > \bar{x})$ \\
Count Below Mean & Number of values below the mean & $N_{\text{below}} = \sum_{i=1}^n (x_i < \bar{x})$ \\
Day of Year of Maximum Value & Day of the year when the maximum value occurs in series & --- \\
Day of Year of Minimum Value & Day of the year when the minimum value occurs in series & --- \\
Kurtosis & Measure of the tailedness of the time series distribution & $G_2 = \frac{\mu_4}{\sigma^4} - 3$ \\
Linear Time Trend & Linear trend coefficient estimated over the entire time series & $b = \frac{\sum_{i=1}^n (t_i - \bar{t})(x_i - \bar{x})}{\sum_{i=1}^n (t_i - \bar{t})^2}$ \\
Longest Strike Above Mean & Longest consecutive sequence of values above the mean & --- \\
Longest Strike Below Mean & Longest consecutive sequence of values below the mean & --- \\
Maximum & Maximum value of the time series & $x_{\text{max}}$ \\
Mean & Mean value of the time series & $\bar{x} = \frac{1}{n}\sum_{i=1}^n x_i$ \\
Mean Absolute Change & Mean of absolute differences between consecutive values & $\frac{1}{n-1} \sum_{i=1}^{n-1} | x_{i+1} - x_{i}|$ \\
Mean Change & Mean of the differences between consecutive values & $ \frac{1}{n-1} \sum_{i=1}^{n-1}  (x_{i+1} - x_{i}) $ \\
Mean Second Derivative Central & measure of acceleration of changes in a time series data & $\frac{1}{2(n-2)} \sum_{i=1}^{n-2}  \frac{1}{2} (x_{i+2} - 2 \cdot x_{i+1} + x_i)
$ \\
Median & Median value of the time series & $\tilde{x}$ \\
Minimum & Minimum value of the time series & $x_{\text{min}}$ \\
Quantile (q = 0.05, 0.95) & Values representing the specified quantiles (5th and 95th percentiles) & $Q_{0.05}, Q_{0.95}$ \\
Ratio Beyond r Sigma (r=1,2,3) & Proportion of values beyond r standard deviations from the mean & $P_r = \frac{1}{n}\sum_{i=1}^{n} (|x_i - \bar{x}| > r\sigma)$ \\
Skewness & Measure of the asymmetry of the time series distribution & $\frac{n}{(n-1)(n-2)} \sum \left(\frac{x_i - \bar{x}}{\sigma}\right)^3$ \\
Standard Deviation & Standard deviation of the time series & $\sqrt{\frac{1}{n}\sum_{i=1}^{n} (x_i - \bar{x})^2}$ \\
Sum Values & Sum of all values in the time series & $S = \sum_{i=1}^{n} x_i$ \\
Symmetry Looking & Measures the similarity of the time series when flipped horizontally & $| \bar{x}-\tilde{x} | < r \cdot (x_{\text{max}} - x_{\text{min}} ) $ \\
Time Series Complexity (CID CE) & measure of number of peaks and valleys & $\sqrt{ \sum_{i=1}^{n-1} ( x_{i+1} - x_{i})^2 }$\\
Variance & Variance of the time series & $\sigma^2 = \frac{1}{n}\sum_{i=1}^{n} (x_i - \bar{x})^2$ \\
Variance Larger than Standard Deviation & check if variance is larger than standard deviation & $\sigma^2 > 1$ \\
\hline
\end{longtable}

\hypertarget{s-ii.-computational-complexity-of-deep-learning-baselines}{%
\section{S-II. Computational Complexity of Deep-Learning
Baselines}\label{s-ii.-computational-complexity-of-deep-learning-baselines}}

This section expands the efficiency comparison in the main text
(Methods, ``Computational Efficiency''), where the LightGBM ``lite
learning'' model trains in 10--15 minutes on a CPU-only laptop. By
contrast, the deep-learning architectures commonly applied to satellite
image time series carry substantially higher training-time and memory
costs. As in the main text, \(n\) denotes the number of training
samples.

Recurrent models such as LSTMs, frequently used for temporal land cover
analysis, exhibit complexity:

\[T_{LSTM} = O(e \cdot n \cdot t \cdot (h^2 + h \cdot i))\]

where \(e\) is the number of epochs, \(t\) the sequence length, \(h\)
the hidden dimensionality, and \(i\) the input feature dimension. The
quadratic dependence on \(h\) and the inherently sequential nature of
LSTMs prevent efficient parallelization across the temporal axis, often
extending training to several hours.

Spatial architectures further escalate these demands. A 2D U-Net scales
linearly with the number of input channels (treating time steps as
channels), but a 3D U-Net, which treats the temporal dimension as a
third spatial axis, follows:

\[T_{3DUNet} = O\left(e \cdot n \cdot \sum_{l} k_l^3 \cdot c_{l-1} \cdot c_l \cdot s_l\right)\]

where the sum runs over network layers \(l\) with kernel size \(k_l\),
input and output channel counts \(c_{l-1}\) and \(c_l\), and spatial
size \(s_l\). The cubic kernel term (\(k_l^3\)) means a
\(3 \times 3 \times 3\) convolution requires 27 operations per voxel,
compared to 9 for a 2D equivalent. Furthermore, the memory (VRAM)
required to store 4D tensors (Space \(\times\) Time \(\times\) Channels)
for large-scale mapping routinely necessitates high-end GPU clusters.

\hypertarget{references}{%
\section*{References}\label{references}}
\addcontentsline{toc}{section}{References}

\hypertarget{refs}{}
\begin{CSLReferences}{1}{0}
\leavevmode\vadjust pre{\hypertarget{ref-xr_fresh_2021}{}}%
Mann, Michael L. 2024. {``{xr\_fresh: Python Package for Feature
Extraction from Raster Data Time Series}.''}
\url{https://doi.org/10.5281/zenodo.12701466}.

\end{CSLReferences}

\end{document}
